% p.35 — الباب الثامن العشر في أحبّ يحبّ (مضاعف في أفعل) — تتمة
\mspage{35}
% Layout: six columns of continuous paradigms (see p.31 note); boxed headings
% (title + uncontracted base form) are given as two consecutive rows. The
% running header is blank/faded in the scan; faint lines below the table are
% show-through from the verso and are not transcribed.
\begin{sarf}{|c|c|c|c|c|c|}
لَمْ تُحِبَّا & لَمْ اُحَبَّ & لَنْ يُحَبَّا & لَتُحِبَّنَّ & لَتُحْبَبْنَانِّ & \thead{أمر حاضر بالنون}\\
لَمْ تُحِبُّوْا & لَمْ نُحَبَّ & لَنْ يُحَبُّوْا & لَتُحِبَّانِّ & لَاُحَبَّنَّ & اَحْبِبَنَّ\\
لَمْ تُحِبِّيْ & \thead{مضارع بلن} & لَنْ تُحَبَّ & لَتُحِبُّنَّ & لَنُحَبَّنَّ & اَحِبَّنَّ\\
لَمْ تُحِبَّا & لَنْ يُحْبِبَ & لَنْ تُحَبَّا & لَتُحِبِّنَّ & \thead{منفي بالنون} & اَحِبَّانِّ\\
لَمْ تُحْبِبْنَ & لَنْ يُحِبَّ & لَنْ يُحْبَبْنَ & لَتُحِبَّانِّ & لَا يُحْبِبَنَّ & اَحِبُّنَّ\\
لَمْ اُحِبَّ & لَنْ يُحِبَّا & لَنْ تُحَبَّ & لَتُحْبِبْنَانِّ & لَا يُحِبَّنَّ & اَحِبِّنَّ\\
لَمْ نُحِبَّ & لَنْ يُحِبُّوْا & لَنْ تُحَبَّا & لَاُحِبَّنَّ & \thead{مجهول بالنون} & اَحِبَّانِّ\\
\thead{مجهول بلم} & لَنْ تُحِبَّ & لَنْ تُحَبُّوْا & لَنُحِبَّنَّ & لَا يُحْبَبَنَّ & اَحْبِبْنَانِّ\\
لَمْ يُحْبَبْ & لَنْ تُحِبَّا & لَنْ تُحَبِّيْ & \thead{مجهول بالنون} & لَا يُحَبَّنَّ & \thead{نهي حاضر بالنون}\\
لَمْ يُحَبَّ & لَنْ يُحْبِبْنَ & لَنْ تُحَبَّا & لَيُحْبَبَنَّ & \thead{أمر غائب بالنون} & لَا تُحْبِبَنَّ\\
لَمْ يُحَبَّا & لَنْ تُحِبَّ & لَنْ تُحْبَبْنَ & لَيُحَبَّنَّ & لِيُحْبِبَنَّ & لَا تُحِبَّنَّ\\
لَمْ يُحَبُّوْا & لَنْ تُحِبَّا & لَنْ اُحَبَّ & لَيُحَبَّانِّ & لِيُحِبَّنَّ & لَا تُحِبَّانِّ\\
لَمْ تُحَبَّ & لَنْ تُحِبُّوْا & لَنْ نُحَبَّ & لَيُحَبُّنَّ & \thead{مجهول بالنون} & لَا تُحِبُّنَّ\\
لَمْ تُحَبَّا & لَنْ تُحِبِّيْ & \thead{مضارع بالنون} & لَتُحَبَّنَّ & لِيُحْبَبَنَّ & لَا تُحِبِّنَّ\\
لَمْ يُحْبَبْنَ & لَنْ تُحِبَّا & لَيُحْبِبَنَّ & لَتُحَبَّانِّ & لِيُحَبَّنَّ & لَا تُحِبَّانِّ\\
لَمْ تُحَبَّ & لَنْ تُحْبِبْنَ & لَيُحِبَّنَّ & لَيُحْبَبْنَانِّ & \thead{نهي غائب بالنون} & لَا تُحْبِبْنَانِّ\\
لَمْ تُحَبَّا & لَنْ اُحِبَّ & لَيُحِبَّانِّ & لَتُحَبَّنَّ & لَا يُحْبِبَنَّ & \thead{تمرين}\\
لَمْ تُحَبُّوْا & لَنْ نُحِبَّ & لَيُحِبُّنَّ & لَتُحَبَّانِّ & لَا يُحِبَّنَّ & اَمَدَّ\\
لَمْ تُحَبِّيْ & \thead{مجهول بلن} & لَتُحِبَّنَّ & لَتُحَبُّنَّ & \thead{مجهول بالنون} & اَدَلَّ\\
لَمْ تُحَبَّا & لَنْ يُحْبَبَ & لَتُحِبَّانِّ & لَتُحَبِّنَّ & لَا يُحْبَبَنَّ & اَكَدَّ\\
لَمْ تُحْبَبْنَ & لَنْ يُحَبَّ & لَيُحْبِبْنَانِّ & لَتُحَبَّانِّ & لَا يُحَبَّنَّ & اَمَلَّ\\
\end{sarf}
